% Original analytic teaching diagram, 2026-11-12.
% Omega = {1,2,3}, F = {empty, {1}, {2,3}, Omega}, P = point mass at 1.
% (a) the four events of F with their probabilities; (b) the completion adds the subsets {2}, {3}, {1,2}, {1,3}
% of the null set {2,3} (and their unions with {1}); probabilities of old events are unchanged.
\definecolor{VB}{HTML}{7188A3}\definecolor{VBF}{HTML}{E5EAF0}
\definecolor{VG}{HTML}{829C8B}\definecolor{VGF}{HTML}{E8EEE9}
\definecolor{VY}{HTML}{B4A16F}\definecolor{VYF}{HTML}{F1ECDD}
\definecolor{VR}{HTML}{A7787D}\definecolor{VRF}{HTML}{F0E5E6}
\definecolor{VP}{HTML}{9588A5}\definecolor{VPF}{HTML}{ECE8F0}
\definecolor{VI}{HTML}{3B4252}\definecolor{VM}{HTML}{6B7280}\definecolor{VL}{HTML}{4C4D4F}
\begin{tikzpicture}[x=1mm,y=1mm,>=Stealth,font=\figurefont\fontsize{8}{10}\selectfont,
lab/.style={text=VM,font=\figurefont\fontsize{7.5}{9}\selectfont}]
% \chip{x}{y}{in1}{in2}{in3}{frame colour}{probability}
\newcommand{\chip}[7]{
  \draw[#6,line width=.8pt,rounded corners=1.2pt] (#1,#2) rectangle ({#1+11},{#2+6});
  \foreach \p/\v in {1/#3,2/#4,3/#5} {
    \ifnum\v=1 \fill[VB] ({#1+\p*2.75},{#2+3}) circle (.95);
    \else \draw[VL,line width=.45pt] ({#1+\p*2.75},{#2+3}) circle (.95); \fi
  }
  \node[text=VI,font=\figurefont\fontsize{7.5}{9}\selectfont] at ({#1+5.5},{#2-2.6}) {$#7$};
}
% row (a)
\node[anchor=east,text=VI] at (-2,24) {$\mathcal{F}$};
\chip{0}{21}{0}{0}{0}{VL}{0}
\chip{12.5}{21}{1}{0}{0}{VL}{1}
\chip{25}{21}{0}{1}{1}{VL}{0}
\chip{37.5}{21}{1}{1}{1}{VL}{1}
\node[text=VI] at (24,12) {(a)};
% row (b)
\begin{scope}[shift={(0,-22)}]
\node[anchor=east,text=VI] at (-2,24) {$\overline{\mathcal{F}}$};
\chip{0}{21}{0}{0}{0}{VL}{0}
\chip{12.5}{21}{1}{0}{0}{VL}{1}
\chip{25}{21}{0}{1}{1}{VL}{0}
\chip{37.5}{21}{1}{1}{1}{VL}{1}
\chip{50}{21}{0}{1}{0}{VY}{0}
\chip{62.5}{21}{0}{0}{1}{VY}{0}
\chip{75}{21}{1}{1}{0}{VY}{1}
\chip{87.5}{21}{1}{0}{1}{VY}{1}
\node[text=VI] at (49,12) {(b)};
\end{scope}
\node[lab,anchor=west] at (52,25) {点自左向右为$1,2,3$};
\end{tikzpicture}
